\subsection{RQ1: Are firms converging towards monoculture, and what is the source of this? }
\paragraph{Firm Similarities} We start from the ELO data per-firm from the leaderboard, which gives per-category ELO ratings, $\theta_{i,k}$ for each firm $i$ and category $k$.
The key challenge is that raw ELO values are in the range $650$ to $1560$,
leading to cosine similarity between most pairs of firms near-one.
This leads to a measurement problem of how to rescale ELOs to understand similarities at a point.
To get around this, we instead track residuals, $r_{i,k} = \theta_{i,k} - \bar{\theta_i}-\bar{\theta_k}$, which subtracts out the per-firm and per-category mean to compute residuals for each category cross firm.
Using this, we can study two things: 1) what is the distribution of pairwise residual similarities, $\cos(\mathbf{r}_i,\mathbf{r_{i'}})$; are there groups of firms that cluster together in this, and 2) how consistent is $\mathbf{r}_i$ over time; are per-firm advantages consistent.
We note that computing $H$ using $\mathbf{r}_{i}$ makes little sense, because $H$ will tend towards $-1/(N-1)$ regardless of the distribution of $\btheta_i$, just due to properties of residuals.

\begin{figure*}
    \includegraphics[width=0.57\textwidth]{figures/residual_cosine_kde.pdf}
    \includegraphics[width=0.43\textwidth]{figures/residual_cosine_heatmap.pdf}
    \caption{\textbf{Left:} distribution of pairwise cosine similarities of the two-way (firm, category) residual $\mathbf{r}_i$, pooled across 133 weekly snapshots (Jan.\ 2025--Jun.\ 2026) for the top 23 labs, against the permutation null that shuffles each category's residuals independently across firms. \textbf{Right:} pairwise residual cosine similarity at the most recent snapshot, with rows/columns reordered by hierarchical clustering to expose community structure.}
    \label{fig:residual_cosine_kde}
\end{figure*}

In Figure~\ref{fig:residual_cosine_kde} we plot both the distribution of pairwise similarities in residuals, along with a heatmap of the similarities for different pairs of firms.
We find that the distribution of residuals is statistically indistinguishable in location from the permutation null: the pooled mean pairwise cosine is $-0.053$ against a null mean of $-0.056$ (Mann--Whitney $p=0.106$; Welch $t$-test on the means, $p=0.366$).\footnote{A Kolmogorov--Smirnov test over the full distributions does reject exact equality ($D=0.068$, $p<10^{-68}$), but with $n=18{,}029$ observed and $360{,}580$ null draws the test has power to detect a modest excess in the tails of the observed distribution -- consistent with the community structure documented below, rather than a shift in central tendency.}
This demonstrates that there is no consistent pattern in the residuals themselves; firm residuals are near-random.
However, the heatmap on the right of Figure~\ref{fig:residual_cosine_kde} demonstrates community structure for which sets of firms are more similar to each other.
We find four major communities. The largest, spanning Alibaba, Baidu, DeepSeek, Google, OpenAI, and xAI, is lopsided toward Math and Creative Writing and away from Coding and English -- consistent with these labs converging on similar large-scale RL-for-reasoning post-training recipes over the same period (e.g.\ o1/o3, Gemini ``thinking,'' DeepSeek-R1, Grok reasoning, Qwen3-thinking), a shared \emph{era}, not necessarily a shared \emph{ability level}, since the firm-level mean is already netted out.
A second community (Meta, Mistral, Cohere, ByteDance, Z.ai) shows close to the mirror-image profile, favoring Creative Writing and Multi-Turn over Math.
A third, smaller community (Anthropic, Moonshot, Xiaomi) shares an Instruction-Following-favoring, English-and-Multi-Turn-avoiding profile; we do not have a ready causal story for this grouping and flag it as a small-$n$ curiosity rather than a finding.
The remaining nine labs (including AI2, NVIDIA, and Amazon) form a looser catch-all cluster leaning toward Coding and English.
This demonstrates that while the overall distribution of firm similarities exhibits little structure, there is structure, with different firms clustering on different patterns.

\begin{figure*}
    \includegraphics[width=\textwidth]{figures/residual_timeseries.pdf}
    \caption{Monthly-aggregated two-way residual $r_{i,k}$ for 12 labs $\times$ 8 categories (96 series) over time, colored by each series' own typical magnitude. Two series are highlighted: Anthropic's Coding residual (persistently positive) and Cohere's Math residual (persistently negative).}
    \label{fig:residual_distribution}
\end{figure*}

Next, we understand how residuals change over time.
In Figure~\ref{fig:residual_distribution}, we plot the change over time in residuals for different firm-category pairs.
While most firms have residuals near-0 and fluctuating, some firms do have persistent trends over time.
For example, both Anthropic in coding and Cohere in math have consistent patterns for residuals, with Anthropic performing better than expected and Cohere worse.
These patterns are significant: a Newey--West (HAC, 12 lags) test of whether the time-mean of the residual differs from zero -- robust to the autocorrelation induced by rating the same flagship model repeatedly between releases -- gives Anthropic/Coding a mean residual of $+11.9$ points ($t=+3.94$, $p=0.00008$) and Cohere/Math a mean of $-28.3$ points ($t=-17.34$, $p<10^{-4}$); both survive Benjamini--Hochberg correction across all 96 firm--category pairs tested ($q<0.001$). More broadly, 60 of the 96 pairs are significant at raw $p<0.05$ (58 after FDR correction), so persistent advantage is not confined to these two examples, though few pairs are as extreme as Cohere/Math or Anthropic/English ($t=-20.18$).
These plots show that, while most firms have little sustained comparative advantage, some firms do feature this, which gives evidence to some differentiation among production functions between firms.

\paragraph{Capability Variation}
We next ask how firms change their capabilities over time by investigating changes in $\Delta \theta_{i,k}$ between timesteps.
The entry of new models in a firms lineup allows us to investigate which capabilities the models most improved upon, and those which improved less.
Here, we restrict our attention to frontier firms; those with at least 3 flagship-changeover events with a mature ($\geq$56-day post-release) rating on both sides of the changeover, so that cluster-robust standard errors by event are not degenerate (this drops Meta, which has only 2 such events in our sample).
We first ask how $\Delta \theta_{i,k}$ relies on the gap, $g_{i,k}$, where $g_{i,k} = \theta_{i,k}-\theta_{i^*,k}$, where $\theta_{i^*,k}$ is the leading firm in category $k$.
Specialized firms should exhibit smaller responses to categories where they are behind, while under monoculture, firms chase a general capability profile, and so respond strongly to gaps.

\begin{figure*}
    \includegraphics[width=0.5\textwidth]{figures/evalchasing_scatter.pdf}
    \includegraphics[width=0.5\textwidth]{figures/evalchasing_forest.pdf}
    \caption{\textbf{Left:} firm--category-fixed-effect-residualized gap $g_{i,k}$ vs.\ subsequent capability change $\Delta\theta_{i,k}$ at flagship changeover, one fitted line per firm (65 events, 520 firm--category observations, 2024--2026). \textbf{Right:} per-firm $\hat\beta$ with 95\% confidence intervals (cluster-robust by event), sorted high to low.}
    \label{fig:evalchasing_scatter}
\end{figure*}

In Figure~\ref{fig:evalchasing_scatter}, we plot the gaps between different models in different categories, along with the response in $\theta_{i,k}$.
We compute $\hat{\beta}$ values for each company, and find that they are all positive and above $\hat{\beta} > 0.4$ (OpenAI $1.06$, xAI $0.91$, Alibaba $0.88$, Anthropic $0.68$, Mistral $0.58$, Google $0.49$, DeepSeek $0.47$).
Moreover, all firms feature a positive confidence interval, even after correcting for noise: because the outgoing flagship's own sampling noise $\epsilon_{\text{old}}$ enters both the gap $g_{i,k}$ (as $-\epsilon_{\text{old}}$) and the subsequent change $\Delta\theta_{i,k}$ (also as $-\epsilon_{\text{old}}$, since it is differenced away between the old and new rating), a naive regression is mechanically biased toward finding $\hat\beta>0$ purely from measurement error. We correct for this using each observation's reported Bradley--Terry sampling variance, subtracting the implied spurious covariance and variance from the moment conditions before solving for $\hat\beta$ (an errors-in-variables correction); the corrected estimates (OpenAI $1.28$, xAI $1.08$, Alibaba $1.21$, Anthropic $0.69$, Mistral $0.75$, Google $0.48$, DeepSeek $0.48$) remain positive and, if anything, larger than the naive ones.
These results demonstrate that gap-chasing is a real phenomenon or strategy implemented by models; rather than lag behind and specialize, they tend to pursue the top models and close the gap across categories.

\begin{table*}
\caption{Robustness of the gap-chasing slope $\hat\beta$ to progressively richer controls. Column (A) is a pooled regression of $\Delta\theta_{i,k,t+1}$ on $\mathrm{Gap}_{i,k,t}$ alone. (B) adds separate firm, category, and time (quarter) fixed effects. (C) replaces the separate firm and category effects with their interaction (firm$\times$category) while keeping time fixed effects. (D), our preferred specification, replaces the time fixed effects with category$\times$time fixed effects, so that common shocks to a category's difficulty in a given quarter (e.g.\ a burst of reasoning-focused post-training) cannot masquerade as gap-chasing. (E) adds $\log$(days since the firm's previous flagship release) as a control -- the available proxy for ``scaling,'' since public parameter/compute figures do not exist for these closed frontier flagships. (F) further applies the same noise correction as Figure~\ref{fig:evalchasing_scatter}. Standard errors (cluster-robust by event) in parentheses. $^{***}p<0.01$.}
\label{tab:beta_robustness}
\begin{tabular}{lcccccc}
\toprule
 & (A) Basic & (B) Model 1 & (C) Model 2 & (D) Model 3 (main) & (E) + Scaling & (F) Full \\
\midrule
$\hat\beta$ (Gap) & +0.354$^{***}$ & +0.314$^{***}$ & +0.334$^{***}$ & +0.338$^{***}$ & +0.354$^{***}$ & +0.341$^{***}$ \\
 & (0.060) & (0.075) & (0.079) & (0.081) & (0.094) & (0.094) \\
\midrule
Firm FE & -- & \checkmark & -- & -- & -- & -- \\
Category FE & -- & \checkmark & -- & -- & -- & -- \\
Time (quarter) FE & -- & \checkmark & \checkmark & -- & -- & -- \\
Firm$\times$Category FE & -- & -- & \checkmark & \checkmark & \checkmark & \checkmark \\
Category$\times$Time FE & -- & -- & -- & \checkmark & \checkmark & \checkmark \\
Noise correction & -- & -- & -- & -- & -- & \checkmark \\
\midrule
$N$ & 520 & 520 & 520 & 520 & 456 & 456 \\
Events & 65 & 65 & 65 & 65 & 57 & 57 \\
\bottomrule
\end{tabular}
\end{table*}

We ask whether such results hold even when varying the underlying model; is the presence of $\hat{\beta}$ merely due to an underspecified model.
In Table~\ref{tab:beta_robustness}, we compare different models and the resulting $\hat{\beta}$.
In particular we compare a basic pooled regression with no controls (A); one with separate firm, category, and quarter fixed effects (B); one that replaces the separate firm and category terms with their interaction while keeping a time effect (C); our preferred specification with firm$\times$category and category$\times$time fixed effects (D); a version that further controls for release cadence via $\log$(days since previous release) (E), since firm size/compute data does not exist for these closed models; and a final version that also applies the sampling-noise correction described above (F).
We find that $\hat{\beta}$ is persistent, being at least $0.3$ across different models.
These results demonstrate that even accounting for other confounds, gap-chasing is a real phenomenon.

Digging more into the full model, we find that a noise corrected model with category and time fixed effects reveals that $\hat\beta$ is essentially unchanged (0.338 vs.\ 0.354 in the basic specification) once we absorb both common per-category shocks over time (category$\times$time, which would otherwise let a shared trend like ``everyone's math score jumps in a given quarter'' masquerade as gap-chasing) and each firm's own persistent comparative advantage in a category (firm$\times$category, which absorbs exactly the kind of effect documented in Figure~\ref{fig:residual_distribution} -- e.g.\ Anthropic's standing edge in Coding). The fact that these fixed effects are doing real work -- and yet leave $\hat\beta$ nearly untouched -- indicates that gap-chasing is a distinct phenomenon from either persistent firm-category advantage or common category-level trends, rather than an artifact of either.
For example, these fixed effects get at the persistent effect that different companies have in some categories.

\paragraph{Real-World Usage}
We next investigate how changes in model performance correspond to real-world usage through an analysis of the OpenRouter dataset.
The OpenRouter dataset features information on the top-50 models used each day, along with their usage, allowing us to estimate how the distribution of models used has changed over time. We match OpenRouter's model slugs to Arena model names by normalizing both to a canonical token representation (stripping provider prefixes and version-string decorations) and, where that fails, applying a small set of hand-verified aliases for renamed or reformatted releases; this matches 192 of 347 distinct OpenRouter models to an Arena name, covering an average of 79\% of monthly token volume (ranging from 61\% to 92\% across Jan.\ 2025--Jun.\ 2026 as new, not-yet-matched releases enter and leave the top of the usage distribution). Unmatched volume (stealth-codenamed releases, embedding models, and not-yet-aliased new checkpoints) is always treated as \emph{not} within any threshold of the top model, so the shares we report below are a conservative floor on the true concentration of usage near the frontier.

\begin{figure*}
    \includegraphics[width=0.5\textwidth]{figures/openrouter_share_by_threshold.pdf}
    \includegraphics[width=0.5\textwidth]{figures/near_top_single.pdf}
    \caption{\textbf{Left:} average share of OpenRouter token volume served by models within $x\%$ of that month's top Arena ``overall'' rating, as $x$ varies; the dashed gray line marks the average matched-token coverage (79\%), which the curve cannot exceed by construction. \textbf{Right:} the average number of models (across the 8 CORE categories) within 10\% of that category's top rating, over time.}
    \label{fig:openrouter_effective_rank}
\end{figure*}

In Figure~\ref{fig:openrouter_effective_rank}, we plot both the fraction of models used by Arena threshold $x$, and the rise in the number of models that are within the top-10\%.
We find that most models used are within 10\% of the best model on LLM Arena, where best is determined per-category, and averaged across categories: the token-share curve rises steeply from 4\% of volume at $x=1\%$ to 70\% at $x=10\%$, then flattens to its 79\% coverage ceiling by $x=20\%$ (78\%, 78\%, and 78\% at $x=20,30,50$) -- an elbow (maximum distance from the chord joining the first and last points of the curve) located exactly at $x=10\%$, which we use as the threshold for the right-hand panel. Moreover, we find that the effective rank of OpenRouter deployment (the inverse Herfindahl index of token share across models) has increased over time; from January 2025 to June 2026, the effective number of models served rose from 7.2 to 22.3 (linear trend $+13.0$/yr, $r=0.85$).
We can map this to the increase in the number of models within the top-10\% in the same time period.
From January 2025 to June 2026, the number of models within the top-10\% has increased from an average of 41.4 to 131.6 across the 8 CORE categories -- roughly a 3.2$\times$ increase.
As a result, there's a greater variation in the number of available choices, leading to diversification in usage.
This partially originates from monoculture; more models with similar profiles leads to larger substitutability, and a more diverse set of models used in practice.
Together, these results suggest that the rise in deployment diversity is not simply usage fragmenting across an unchanged pool of top models, but a genuine broadening of the near-frontier itself -- consistent with the view that eval-chasing (documented above) produces more, not fewer, viable substitutes at the top of the market.

\subsection{RQ2: Based on evaluation data, can we recover any information on the similarity of production functions between firms? }
\paragraph{Benchmark Scaling Curves}
We investigate what plausible production functions look like by looking at the ways in which open source models scale.
Using information from the Open LLM Leaderboard's per-model results, joined against each model's parameter count, restricted to 6 objective capability dimensions (IFEval, BBH, MATH, GPQA, MuSR, MMLU-Pro) and to labs with a genuine scaling ladder to fit against (at least 7 models spanning at least a $10^{0.6}\approx4\times$ range in parameter count, with at least 4 distinct sizes) -- 15 labs and 278 models survive this screen.
From this, we can plot how the number of parameters impacts model performance and scaling.

\begin{figure*}
    \includegraphics[width=\textwidth]{figures/scaling_curves_normalized.pdf}
    \includegraphics[width=\textwidth]{figures/scaling_beta_heatmap.pdf}
    \caption{\textbf{Top:} score vs.\ $\log_{10}$(parameters, B) for each of the 6 objective categories, standardized ($z$-score) within category so slopes are comparable across panels; gray points/lines are individual models/firm-level fits, the four most data-rich labs are highlighted, and the bold black line is the pooled (all-labs) fit. \textbf{Bottom:} the fitted slope $\hat\beta_{i,k}$ for each of the 15 labs (columns, sorted by mean slope) and 6 categories (rows).}
    \label{fig:scaling_curves_normalized}
\end{figure*}

In Figure~\ref{fig:scaling_curves_normalized}, we plot how open source models scale as a function of the number of parameters across different categories.
Scores within each category are standardized to zero mean and unit variance across the whole pool of 278 models before fitting, so that a slope is directly comparable across categories despite the underlying benchmarks having very different raw scales and ranges (e.g.\ MATH scores span a much wider raw range than IFEval); without this, a shared $y$-axis across the six panels would be meaningless.
We find slight heterogeneity across categories and firms, but generally find the same qualitative pattern: a positive, concave relationship between $\log_{10}$(parameters) and score in every one of the 6 categories, with the pooled slope ranging from $+0.75$ (IFEval, the flattest) to $+1.15$ (MMLU-Pro, the steepest), a $\sim$1.5$\times$ spread that is real but modest (95\% bootstrap CIs on adjacent categories' slopes overlap only at the extremes).
On the right, we plot the slope for different firms and categories.
Decomposing the variance of the fitted slope matrix $\hat\beta_{i,k}$ into firm, category, and firm$\times$category components, we find \emph{more} differentiation across firms than across categories in raw terms: the row (firm) means account for 70\% of the decomposed variance, column (category) means for 7\%, and the firm$\times$category interaction for the remaining 22\%. This is the opposite emphasis from the held-out prediction result in Table~\ref{tab:scaling_function_fits} below, where firm identity \emph{alone} buys almost nothing (0.8\%) -- the reconciliation is that per-firm average slopes are estimated noisily (many labs have few models or a narrow size range), so a firm's average slope can look very different from another's in-sample without that difference generalizing; it is specifically the firm$\times$category cell, not either main effect alone, that carries reliable, out-of-sample signal. AI2 and Arcee AI have the highest average slope (both $\approx1.9$), while TII and EleutherAI have the lowest (both statistically indistinguishable from zero, at $-0.09$ and $0.01$ respectively).
This suggests that a purely firm-level notion of ``scales better'' is not by itself a reliable summary of the data; scaling advantages are substantially category-specific, which is precisely the kind of heterogeneity that a scalar per-firm production function cannot represent.


\paragraph{Possible $F_i$ and elimination.} Show how we can eliminate some $F_i$ and give certain properties for others.
\begin{table}[t]
    \centering
    \begin{tabular}{lccc}
        \toprule
        Model for $F$
        & Held-out error $\downarrow$
        & Firm gain
        & Category gain \\
        \midrule
        Linear
        & 17.626 & -- & -- \\
        Log
        & 17.039 & -- & -- \\
        Power law
        & 17.002 & -- & -- \\
        Power law + firm
        & 16.862 & 0.8\% & -- \\
        Power law + firm $\times$ category
        & \textbf{11.507} & 0.8\% & 31.8\% \\
        \bottomrule
    \end{tabular}
    \caption{5-fold cross-validated held-out error (RMSE, split by model so a model's scores never straddle train/test) for progressively richer functional forms $F$ relating $\log_{10}$(params) to standardized score. ``Firm gain'' is the \% reduction in held-out error from adding firm-specific $(a_i,b_i)$ to the pooled power-law fit (curvature $\hat\gamma\approx0.13$, shared across firms); ``Category gain'' is the further \% reduction from letting $(a,b)$ vary by firm$\times$category instead of firm alone. Functional form barely matters (rows 1--3 are within noise of each other); nearly all of the real heterogeneity in this data is firm$\times$category-specific, not firm-specific.}
    \label{tab:scaling_function_fits}
\end{table}

We next discuss how the different patterns we have seen so far correspond to different families of production functions.
We present five properties that we have observed at least partially empirically that might be replicable through a production function.
These properties are: (i) diminishing returns to scale (Table~\ref{tab:scaling_function_fits}'s fitted power-law exponent $\hat\gamma\approx0.13\ll1$); (ii) category heterogeneity in scaling response (the $0.75$--$1.15$ slope spread by category in Figure~\ref{fig:scaling_curves_normalized}, and the $31.8\%$ held-out gain from category-specific parameters in Table~\ref{tab:scaling_function_fits}); (iii) firm heterogeneity (the same figure's firm-level slope spread, AI2/Arcee AI vs.\ TII/EleutherAI); (iv) persistent, category-specific comparative advantage (Anthropic/Coding, Cohere/Math from RQ1); and (v) the possibility of cross-category spillovers, which none of our empirical results directly test for but which the theory motivates as a further degree of freedom.
We then ask which production functions can preserve these.
We find that some level of flexibility is necessary in production functions; for example, the need to model spillovers.
However, once spillovers are modeled, say with $F_i(\mathbf{x}_i) = A_i \log(\mathbf{x}_i+\mathbf{1})$, this presents sufficient heterogeneity to capture all of these effects.
The key thing is that the matrix $A_i$ needs to be sufficiently flexible, which requires unrestrained $A_i$.

\begin{table*}
    \centering
    \begin{tabular}{lccccc}
\toprule
$F_i$ family & Diminishing & Category & Firm & Persistent comp. & Cross-category \\
             & returns     & heterog. & heterog. & advantage$^\dagger$    & spillovers$^\ddagger$ \\
\midrule
$a_i\log(1+x)\mathbf{1}$                & \checkmark & $\times$   & \checkmark & $\times$   & $\times$ \\
$A_i\log(1+x)$, diagonal $A_i$          & \checkmark & \checkmark & \checkmark & \checkmark & $\times$ \\
$A_i\log(1+x)$, general $A_i$          & \checkmark & \checkmark & \checkmark & \checkmark & \checkmark \\
More flexible $F_i$ (e.g.\ power law)   & \checkmark & \checkmark & \checkmark & \checkmark & \checkmark \\
\bottomrule
\end{tabular}
\caption{Which $F_i$ families are even capable of generating each empirical property. $\dagger$: a scalar $a_i$ cannot produce persistent, category-specific advantage -- directly falsified by RQ1's finding that Anthropic is simultaneously significantly \emph{positive} on Coding ($t=+3.94$) and significantly \emph{negative} on English ($t=-20.18$). $\ddagger$: the one property in this table we have not directly tested; rows 3--4 can represent it, but nothing in our results isolates a genuine cross-category spillover from ordinary per-category heterogeneity (row 2 already explains everything else observed).}
\label{tab:fi_requirements}
\end{table*}


\subsection{RQ3: How would changes in the evaluation procedure impact our results?}
We finally consider a semi-synthetic simulation, where we ask how variation in $\bw$ impacts downstream effects such as $H$ and our ability to identify production functions.
We calibrate the simulation to real data: the initial capability profile $\theta_i^0$ is each frontier firm's real, latest 8-category profile; the release-noise pool is bootstrap-resampled from real per-event residuals (preserving their real cross-category correlation structure, not i.i.d.\ Gaussian); and the baseline production matrix $B_0$ (an isotropic term $\beta_0 I$ plus a shared rank-1 correction) is estimated by ridge regression on the real gap/$\Delta\theta$ events from RQ2. Against this common calibration we simulate three plausible $F_i$ families: (i) \emph{homogeneous}, every firm shares the identical $B_0$; (ii) \emph{parallel, scaled}, every firm shares $B_0$'s direction but is scaled by its own real, previously-estimated $a_i$ (Theorem 9's hypothesis); and (iii) \emph{fully heterogeneous}, every firm gets $B_0$ plus an independent random $K\times K$ perturbation. For the control result, we hold $\bw$ fixed at $\bw(\lambda)=(1-\lambda)\bw^{A}+\lambda \bw^{B}$ for 30 simulated release cycles (a steady state) and measure the resulting real (non-residualized) monoculture $H(\theta^*(\bw(\lambda)))=\overline{\cos}(\theta_i^*,\theta_{i'}^*)$, where $\bw^A,\bw^B$ concentrate 70\% of evaluator weight on Coding and Creative Writing respectively.

\begin{figure}[t]
    \centering
    \includegraphics[width=0.9\columnwidth]{figures/rq3_evaluator_control.pdf}
    \caption{Real monoculture $H(\theta^*(\bw(\lambda)))$ as evaluator weight shifts from Coding ($\lambda=0$) to Creative Writing ($\lambda=1$), for each of the three $F_i$ families, median and 10--90th percentile band over 25 simulation replicates. Under homogeneous technology $H$ is pinned at its ceiling ($\approx1$) for every $\lambda$: with no real technology heterogeneity, the evaluator has no control over monoculture at all. Under the other two families $H$ swings substantially (roughly $0.45$--$0.85$) but non-monotonically in $\lambda$, peaking at intermediate blends rather than at either pure-category extreme.}
    \label{fig:rq3_evaluator_control}
\end{figure}

Figure~\ref{fig:rq3_evaluator_control} shows the resulting control result directly: real technology heterogeneity is necessary for the evaluator to have any leverage over $H$ at all, and even where leverage exists, it need not behave monotonically in the obvious way (a balanced evaluation can induce \emph{more} monoculture than a specialized one).

\begin{table*}
\caption{Response rank and range of induced $H$ for four evaluation designs, all measured against the same fully-heterogeneous $F_i$ draw and the same $T=30$ simulated evaluations (except ``Single evaluation,'' which never varies $\bw$). ``Small/Moderate/Diverse'' differ only in the rank $r\in\{1,2,7\}$ of the subspace of the $(K-1)$-dimensional simplex tangent space that $\bw$'s i.i.d.\ perturbations are drawn from -- the \emph{same} number of evaluations, spent along more or fewer independent directions. $\hat r_i$ is the effective (participation-ratio) rank of each firm's realized, direction-normalized investment-signal sequence $g_t=K\,\bw_t\odot\mathrm{gap}_t$, the empirical analogue of $r_i=\mathrm{rank}(Z_i)$ in the identification proposition.}
\label{tab:semi_synthetic}
    \begin{tabular}{lccccc}
\toprule
Evaluation design & \# distinct $\mathbf{w}$\ & Response rank $\hat r_i$ & Unidentified dims. $K(K-\hat r_i)$ & Range of $H$ induced & $F_i$ identified? \\
\midrule
Single evaluation & 1 & 1.00 & 56.0 & [0.64, 0.65] & No \\
Small perturbations & 30 & 2.47 & 44.2 & [0.54, 0.64] & No \\
Moderate variation & 30 & 2.95 & 40.4 & [0.41, 0.66] & No \\
Diverse variation & 30 & 4.10 & 31.2 & [0.21, 0.61] & No \\
\bottomrule
\end{tabular}
\end{table*}

In Table~\ref{tab:semi_synthetic}, we consider different levels of variation in $\bw$: each row uses the same number ($T=30$) of simulated evaluations, differing only in whether $\bw$'s month-to-month perturbations are confined to a rank-1, rank-2, or rank-7 (of a maximum possible rank $K-1=7$, since $\bw$ lives on the probability simplex) subspace, with total perturbation variance held fixed across rows so that ``richer'' does not just mean ``bigger.''
We find that larger perturbations can lead to larger response ranks, which can reduce the number of unidentified dimensions: $\hat r_i$ rises from $1.00$ (a single, unchanging evaluation) to $2.47$, $2.95$, and $4.10$ as the design's rank increases, and unidentified dimensions fall correspondingly from $56.0$ to $31.2$.
However despite this, we cannot fully identify $F_i$ from the data; there are still many unidentified dimensions: even the most diverse design tested leaves $31.2$ of the $K(K-1)=56$ possible degrees of freedom -- 56\% of the production matrix -- unidentified, and $\hat r_i$ never approaches the full rank $K=8$ that would be required for exact identification. We trace this shortfall to the multiplicative structure of the realized investment signal $g_t=K\,\bw_t\odot\mathrm{gap}_t$: because each firm's competitive gap has its own (largely $\bw$-independent) dynamics, richly-designed variation in $\bw$ does not translate one-for-one into richly-varying realized investment signals, so identification remains partial even when the evaluator does everything a rank-counting argument would suggest it should.
At the same time, varying $\bw$ does change the range of $H$, showing that $\bw$ has an impact on $H$.
Generally, variation in $\bw$ tends to lower $H$, reflecting the theory from Section~\ref{subsec:category_distro} -- more precisely, richer designs widen the achievable range of $H$ almost entirely by pulling its \emph{lower} bound down (from $0.64$ down to $0.21$), while the upper bound stays roughly fixed ($0.61$--$0.66$) across all four designs; the evaluator's leverage is over how much monoculture can be \emph{undone}, not over how much can be added beyond what limited variation already permits.
Taken together, these results suggest that evaluation design is a genuine, usable lever on monoculture even though it is not, in this calibrated environment, a sufficient instrument for fully recovering firms' underlying production technology -- control and identification are both real but partial, and neither comes for free.
